\section{Results}
\label{sec:results}

This section reports speculative-decoding outcomes from the new
RTX4090 CUDA-graph run family summary artifact
(\texttt{RTX4090\_Cuda\_Graph/all\_configs\_summary.csv}) and
regime-matched baseline context from the deterministic and stochastic baseline
CSV artifacts. We emphasise paired comparative metrics
($S$, $\alpha$, $B_{\text{eff}}$) and task deltas to evaluate the
speed--quality tradeoff.

\begin{table*}[t]
\centering
\caption{RTX4090 CUDA-graph run family summary (Qwen2.5-3B target).
$S$ is paired speedup reported by the summary artifact; CI$_S$ denotes paired
bootstrap CI status; $\alpha$ is token-level acceptance; $B_{\text{eff}}$ is
accepted tokens per verification cycle.}
\label{tab:headline}
\small
\begin{tabular}{l c c c c c c c c}
\hline
\textbf{Config} & $S$ & CI$_S$ & $\alpha$ & $B_{\text{eff}}$ & $\Delta$GSM8K & $\Delta$MMLU & $\Delta$CNN/DM & $R_{\text{tok}}$ \\
\hline
baseline, det        & 1.0000 & -- & -- & -- & -- & -- & -- & 11.03 \\
baseline, stoch      & 1.0000 & -- & -- & -- & -- & -- & -- & 10.83 \\
0.5B, $k{=}4$, det   & \textbf{3.0674} & planned & 0.4212 & 1.66 & 0.33 & 0.00 & -0.02 & 34.09 \\
0.5B, $k{=}4$, stoch & 2.9394 & planned & 0.4080 & 1.61 & 1.67 & -0.20 & -0.03 & 32.06 \\
0.5B, $k{=}8$, det   & 2.4542 & planned & 0.3515 & 2.72 & 0.33 & 0.00 & -0.05 & 27.77 \\
0.5B, $k{=}8$, stoch & 2.4063 & planned & 0.3444 & 2.67 & 1.00 & 0.00 & -0.38 & 26.78 \\
0.5B, $k{=}16$, det  & 1.8050 & planned & 0.2539 & 3.76 & 0.33 & 0.00 & -0.01 & 21.90 \\
0.5B, $k{=}16$, stoch& 1.8185 & planned & 0.2516 & 3.73 & 0.34 & -0.20 & -0.19 & 21.51 \\
1.5B, $k{=}4$, det   & 2.8726 & planned & 0.4581 & 1.81 & -0.33 & 0.00 & -0.03 & 31.67 \\
1.5B, $k{=}4$, stoch & 2.7860 & planned & 0.4421 & 1.75 & 1.67 & 0.00 & -0.34 & 30.21 \\
1.5B, $k{=}8$, det   & 2.4066 & planned & 0.3929 & 3.04 & 0.00 & 0.00 & -0.06 & 26.78 \\
1.5B, $k{=}8$, stoch & 2.3444 & planned & 0.3822 & 2.96 & 3.67 & 0.20 & -0.47 & 25.59 \\
1.5B, $k{=}16$, det  & 1.8068 & planned & 0.2968 & 4.39 & 0.67 & 0.00 & -0.03 & 21.49 \\
1.5B, $k{=}16$, stoch& 1.7773 & planned & 0.2928 & 4.36 & 1.00 & -0.20 & -0.36 & 20.70 \\
\hline
\end{tabular}
\end{table*}

Figure~\ref{fig:specdec-vs-baseline-metrics} complements
Table~\ref{tab:headline} with a visual baseline-relative view across all
configurations, including speedup, latency/throughput deltas, token-timing
deltas, and acceptance dynamics.

\begin{figure*}[t]
\centering
\includegraphics[width=0.98\textwidth]{figures/specdec_configs_vs_baseline_metrics.pdf}
\caption{Speculative-decoding metrics for all RTX4090 configurations relative
to regime-matched baseline. Top-left: paired speedup $S$ (baseline at
$S{=}1$). Top-right: mean latency reduction and throughput gain percentages vs
baseline. Bottom-left: TTFT and TPOT percentage deltas vs baseline. Bottom-right:
acceptance dynamics ($\alpha$ and $B_{\text{eff}}$) across the same
configurations.}
\label{fig:specdec-vs-baseline-metrics}
\end{figure*}

\subsection{RQ1: Does speculative decoding accelerate inference?}

All 12 speculative configurations exceed $S{=}1$ in the RTX4090 CUDA-graph
run family, with observed speedups from $1.7773$ to $3.0674$. The best
configuration is 0.5B with $k{=}4$ in deterministic mode. Baseline context
from regime-matched baseline files shows mean latencies of 11.41s
(deterministic) and 11.65s (stochastic) per sample, so the measured
speculative gains are substantial in wall-clock terms for fixed workload and
prompts. Statistical confirmation remains reported via paired uncertainty
analysis on a coherent run family.

\subsection{RQ2: How do draft length and draft size affect acceptance and speed?}

Draft length shows a monotonic decline in speedup over this grid. Averaging
over drafts and regimes, mean speedup is highest at $k{=}4$ ($2.9164$), then
$k{=}8$ ($2.4029$), then $k{=}16$ ($1.8019$). This is consistent with acceptance
dynamics: as $k$ increases, $B_{\text{eff}}$ increases (1.7075 $\rightarrow$
2.8475 $\rightarrow$ 4.0600) but $\alpha$ decreases (0.4324 $\rightarrow$
0.3678 $\rightarrow$ 0.2738), and the extra draft/verify cost dominates net
benefit.

Draft size has a smaller global effect than $k$ in this grid, and the 0.5B
draft is slightly better on mean speedup (mean $S$: 2.4151 for 0.5B versus
2.3323 for 1.5B).

\begin{table*}[t]
\centering
\caption{Representative deterministic operating points showing the non-linear
relationship between acceptance statistics and end-to-end speedup.}
\label{tab:core-bottleneck}
\small
\begin{tabular}{l c c c c l}
\hline
\textbf{Draft} & $k$ & $\alpha$ & $B_{\text{eff}}$ & $S$ & \textbf{Bottleneck interpretation} \\
\hline
0.5B & 4  & 0.4212 & 1.66 & 3.0674 & Sweet spot: draft/verify overhead and acceptance are balanced. \\
0.5B & 16 & 0.2539 & 3.76 & 1.8050 & Drafting overload: acceptance gain cannot cover longer draft loop. \\
1.5B & 4  & 0.4581 & 1.81 & 2.8726 & Heavier draft increases memory traffic despite higher acceptance. \\
1.5B & 16 & 0.2968 & 4.39 & 1.8068 & Diminishing returns: larger accepted block but lower net speedup. \\
\hline
\end{tabular}
\end{table*}

Table~\ref{tab:core-bottleneck} makes the central systems result explicit:
increasing accepted block size does not guarantee higher throughput on this
consumer GPU. As $k$ grows, verification-loop overhead and draft-side memory
pressure rise faster than useful acceptance, so practical policy should
prioritise per-cycle latency minimisation rather than raw acceptance increase.

\subsection{RQ3: Deterministic versus stochastic regime}

Deterministic decoding is faster in 5 of 6 matched draft--$k$ pairs. The
deterministic advantage ranges from $-0.0135$ to $0.1280$ in absolute $S$ and
averages $0.0568$ across the six matched pairs; the single exception is
0.5B, $k{=}16$, where stochastic is slightly faster. This still supports a
deployment preference for deterministic regime when latency is primary and
output diversity is not required.

Quality deltas show different stability by task. CNN/DailyMail stays close to
baseline but is consistently negative (from $-0.47$ to $-0.01$). MMLU
variation is tightly bounded (from $-0.2$ to $0.2$). GSM8K is the most
sensitive axis (from $-0.33$ to $3.67$), especially under stochastic settings.
The regime split is substantial: GSM8K delta spread is 1.00 points in
deterministic mode versus 3.33 points in stochastic mode, and an auxiliary
pairwise output comparison at $k{=}4$ shows much higher stochastic
disagreement with baseline (83.1\%) than deterministic disagreement (9.4\%).
This supports an entropy-sensitive interpretation: high-constraint reasoning
tasks amplify the interaction between sampling noise and speculative
correction, while lower-entropy tasks remain more stable.
Under a conservative quality filter ($\max |\Delta| \le 1$ across tasks),
9 configurations remain, with best speedup at 0.5B, $k{=}4$, deterministic.

While Table~\ref{tab:headline} reports means from the current summary artifact,
camera-ready statistical reporting will pair these means with CI$_S$ and
significance values computed from a single coherent rerun family.

Implementation details for the CUDA-graph-capable runtime setup are described
in Section~\ref{sec:protocol}.

\subsection{Summary of empirical operating point}

The current evidence supports a practical recommendation: use deterministic
speculative decoding with a smaller block ($k{=}4$), where speedup is highest,
and prefer the 0.5B draft for this hardware profile. Higher $k$ increases
accepted block size but lowers net speedup in this setup, indicating that
acceptance gains no longer offset added cycle cost. ACSD results in
Section~\ref{sec:acsd} show that adaptive control can approach these best
fixed-policy speeds while adding explicit rescue/fallback guardrails.
